\documentclass[11pt]{article}
\usepackage[margin=1in]{geometry}
\usepackage[T1]{fontenc}
\usepackage{lmodern,amsmath,amssymb,amsthm,mathtools,booktabs}
\usepackage[hidelinks]{hyperref}
\newcommand{\tr}{\operatorname{tr}}
\newcommand{\ket}[1]{\lvert#1\rangle}
\newcommand{\bra}[1]{\langle#1\rvert}
\newcommand{\N}{\mathcal N}
\newcommand{\T}{\mathcal T}
\newcommand{\R}{\mathcal R}
\newcommand{\C}{\mathcal C}
\newcommand{\F}{F_{\mathrm e}}
\newcommand{\norm}[1]{\lVert#1\rVert}
\newtheorem{theorem}{Theorem}
\newtheorem{proposition}{Proposition}
\newtheorem{lemma}{Lemma}
\title{Exact rank-one encoding and recovery for the Li--Jiang channel}
\author{Haobo Ma \and Wenlin Zhang}
\date{20 October 2026}
\begin{document}
\maketitle
\begin{abstract}
For the Li--Jiang noise family, we prove that the partial-trace decoder can
be chosen in a jointly optimal rank-one encoder--decoder pair when the
logical dimension is any $d\ge2$, for every $0<p<1$, over all complex
rank-one CPTNI encoders and all CPTNI decoders. The qutrit proof combines
Li--Jiang's rank-one replacement bound with a sufficient condition for a
unitary recovery to solve the full recovery SDP; a coarser fidelity estimate
extends the argument to $d\ge4$. In every logical dimension
$d\ge2$, we determine the full complex rank-one optimum with partial trace
fixed, by a two-amplitude reduction and a specified quartic root.
An explicit decoder SDP dual certificate proves that partial trace is an
optimal decoder for every encoder in this real family. We give an exact
algebraic certificate for the joint qutrit optimum at $p=8/11$.
\end{abstract}

\section{The channel and the joint question}
Li and Jiang~\cite{LiJiang} construct a noise family whose unrestricted
optimal fidelity requires higher-rank encoding. We study whether the
partial-trace decoder can be chosen in a jointly optimal rank-one pair,
and determine the exact fidelity when that decoder is fixed.
Theorems~\ref{thm:qubit} and~\ref{thm:joint} answer the
joint question in every dimension. The argument uses a general
recovery criterion and a coarse feasible fidelity above $3/4$; it does not
require a symmetry reduction of the nonconvex encoder set.
Theorem~\ref{thm:fixed} gives the fixed-decoder
formula for all $d$, and Proposition~\ref{prop:decoder} supplies a matching
decoder certificate throughout the attaining family.

We use the model of Eqs.~(18)--(23) of Ref.~\cite{LiJiang}. Let
$L=S=\mathbb C^d$, $C=L\otimes S$, $\tau=I_d/d$, and
\[
 \ket\Omega=d^{-1/2}\sum_{i=0}^{d-1}\ket{ii},\quad
 P=\ket\Omega\bra\Omega,\quad D_i=I_d\otimes\bra i.
\]
For $0<p<1$ set
\begin{gather}
 \lambda_1=1-\frac p{d^2-2},\quad
 \lambda_4=\frac{d^2-1+p}{d^2},\quad\lambda_5=\frac{1-p}{d^2},\notag\\
 \lambda_2=\frac{(1-\lambda_1)(d^2\lambda_1-1)}{d^2\lambda_4},\qquad
 \lambda_3=\frac{2(1-\lambda_1)^2}{d^2\lambda_4},\notag\\
 a=\sqrt{\frac{d\lambda_4}{d^2-1}},\quad b=\sqrt{d\lambda_5},\quad
 Q=a(I-P)+bP.\label{eq:parameters}
\end{gather}
The noise action is
\begin{align}
 \N_p(X)&=\lambda_3\tr(X)I_C+
 (\lambda_1-\lambda_3)\tr_S(QXQ)\otimes I_S
 +\lambda_2\sum_{i,j}E_{ij}XE_{ij}^\dagger,\label{eq:noise}\\
 E_{ij}&=\frac{\ket{ij}\bra\Omega-\sqrt{\lambda_5}D_j^\dagger D_iQ}
 {\sqrt{\lambda_4}}.\notag
\end{align}
A rank-one CPTNI encoder is $\C_V(X)=VXV^\dagger$, with
$V\in\mathbb C^{d^2\times d}$ and $V^\dagger V\preceq I_d$.
For logical Kraus operators $L_k$, fidelity means the unrenormalized quantity
\begin{equation}\label{eq:fidelity}
 \F(\tau,\Phi)=d^{-2}\sum_k|\tr L_k|^2.
\end{equation}
Let $f_d(p)$ be the maximum with decoder $\tr_S$ fixed, and $F_1^{\rm joint}(p,d)$
the maximum allowing arbitrary CPTNI decoders as well.

\section{The decoder reduction and the qubit case}
Define
\[
 e_i=\frac{\ket i\bra\Omega-\sqrt{\lambda_5}D_iQ}{\sqrt{\lambda_4}}.
\]
Since $E_{ij}=(I\otimes\ket j)e_i$, summing over $j$ in
\eqref{eq:noise} gives the exact identity
\begin{equation}\label{eq:factorization}
 \N_p(X)=\T_p(X)\otimes I_S/d,\qquad
 \T_p(X)=d\lambda_3\tr(X)I_d+
 d(\lambda_1-\lambda_3)\tr_S(QXQ)+d\lambda_2\sum_i e_iXe_i^\dagger.
\end{equation}
Thus every CPTP physical decoder $\mathcal D:C\to L$ induces the CPTP
map $\R(X)=\mathcal D(X\otimes I_S/d)$, with
$\mathcal D\N_p=\R\T_p$. Conversely $\R$ is realized by
$\mathcal D=\R\tr_S$. The decoder Choi variable therefore has dimension
$d^2\times d^2$ rather than $d^3\times d^3$.

Trace loss cannot help. Completing a CPTNI decoder to a CPTP decoder adds
nonnegative terms to \eqref{eq:fidelity}. For a fixed decoder, fidelity is a
positive quadratic form $q(V)$. Complete the polar decomposition of a
contraction as $V=UH$, $U^\dagger U=I_d$, $0\preceq H\preceq I_d$.
Then $V_\pm=U(H\pm i\sqrt{I-H^2})$ are isometries, and
\[
 \frac{q(V_+)+q(V_-)}2-q(V)=\frac14q(V_+-V_-)\ge0.
\]
An isometric encoder and CPTP decoder attain the joint maximum. This is the
contraction argument of Lemma~S1 of Ref.~\cite{LiJiang}.

\begin{theorem}\label{thm:qubit}
For every $0<p<1$, $F_1^{\rm joint}(p,2)=f_2(p)$.
In particular, the decoder can be chosen as $\tr_S$ in a globally optimal
joint rank-one pair.
\end{theorem}
\begin{proof}
For every qubit channel $\mathcal A$, a recovery maximizing its average
input--output fidelity can be chosen unitary~\cite[Theorem~1]{KBF}.
Since $F_{\rm av}(\Phi)=(2\F(I_2/2,\Phi)+1)/3$ for a qubit CPTP map,
the same is true for our objective. This theorem applies to
$\mathcal A=\T_p\C_V$ for each isometry $V$. If $M_k$ are Kraus
operators of $\T_p$ and $U$ is such a unitary recovery, cyclicity gives
\[
 \frac14\sum_k|\tr(UM_kV)|^2
 =\frac14\sum_k|\tr(M_kVU)|^2.
\]
The matrix $VU$ is still an isometry. The unitary recovery is therefore
absorbed into the encoder, proving the equality. The reverse inequality
follows from feasibility of $\tr_S$.
\end{proof}
The cited qubit theorem follows by removing the recovery's Bloch translation
without changing its objective and then using that every unital qubit channel
is a mixture of unitary channels~\cite[Appendix~A]{KBF}. Its conclusion does
not extend to arbitrary channels in dimension three.

\subsection{A recovery criterion and joint optimality in every dimension}
For a CPTP channel $\mathcal A:\mathsf L(\mathbb C^d)\to
\mathsf L(\mathbb C^d)$, write
\[
 u(\mathcal A)=\max_{U\ {\rm unitary}}\F(I_d/d,\operatorname{Ad}_U\mathcal A),
 \qquad
 t(\mathcal A)=\max_{\R\ {\rm CPTP}}\F(I_d/d,\R\mathcal A).
\]
The unitary quasi-inverse theorem used above is special to qubits:
higher-dimensional channels can require nonunitary quasi-inverses~\cite{SSMZK}.
The following sufficient condition applies in any dimension.

\begin{lemma}\label{lem:unitary}
If $f=u(\mathcal A)>0$ and
\begin{equation}\label{eq:threshold}
 (d+1)f-(d-1)\sqrt f-1\ge0,
\end{equation}
then $t(\mathcal A)=u(\mathcal A)$. If $d=3$ and
$t(\mathcal A)\ge3/4$, an optimal recovery can be chosen unitary.
\end{lemma}
\begin{proof}
Absorb a maximizing unitary into $\mathcal A$, so identity recovery
maximizes the unitary objective. For Kraus operators $A_i$, stationarity
under $\exp(i\theta H)$ for every Hermitian $H$ implies that
\[
 B=\sum_i\overline{\tr A_i}\,A_i
 \quad\text{is Hermitian}.
\]
Indeed the derivative of the fidelity is $-2d^{-2}\operatorname{Im}\tr(HB)$.
The operator $A_0=B/(d\sqrt f)$ is a normalized linear combination of the
Kraus operators, hence can be included in a Kraus representation and is a
contraction. It is Hermitian and has trace $d\sqrt f$. Since each of its
eigenvalues is at most one,
\begin{equation}\label{eq:mineig}
 \lambda_{\min}(A_0)\ge d\sqrt f-(d-1).
\end{equation}

Use input-first vectorization $\ket{M}\!\rangle=\sum_j\ket j\otimes M\ket j$.
The recovery SDP has objective matrix $K=J(\mathcal A^\dagger)/d^2$,
as in~\eqref{eq:decoder-sdp}. Set
\[
 Y=\frac{\sqrt f}{d}A_0^T,\qquad S=Y\otimes I-K,
 \qquad e=\ket I\!\rangle/\sqrt d,\quad P_\perp=I-ee^\dagger.
\]
Then $\tr Y=f$ and $Se=0$, since
$K\ket I\!\rangle=d^{-2}\ket{B^\dagger}\!\rangle$
and $(Y\otimes I)\ket I\!\rangle=(\sqrt f/d)\ket{A_0}\!\rangle$.
Trace preservation gives $\tr K=1/d$, and
$e^\dagger Ke=f/d$, so
$\tr(P_\perp K P_\perp)=(1-f)/d$.
For every unit vector $v\perp e$, positivity of $K$ and~\eqref{eq:mineig} give
\[
 v^\dagger Sv\ge\lambda_{\min}(Y)-\tr(P_\perp K P_\perp)
 \ge\frac{(d+1)f-(d-1)\sqrt f-1}{d}.
\]
Thus $S\succeq0$ under~\eqref{eq:threshold}; $Y$ and identity recovery
are matching dual and primal solutions.

For the qutrit assertion, take a recovery attaining $t=t(\mathcal A)$.
Cyclicity of Kraus traces gives $\F(I_3/3,\R\mathcal A)=
\F(I_3/3,\mathcal A\R)$. Apply the CPTP version of
Li--Jiang's Theorem~1~\cite{LiJiang} with channel $\mathcal A$ and encoder
$\R$, both acting on $\mathbb C^3$. Its rank-one CPTP replacement is a
square isometry, hence a unitary. With $\varepsilon=1-t\le1/4$ the bound is
\[
 f\ge t-\frac29\bigl(1-\sqrt{1-3\varepsilon}\bigr)^2
 \ge\frac34-\frac1{18}=\frac{25}{36}.
\]
For $\sqrt f\ge5/6$, $4f-2\sqrt f-1\ge1/9>0$.
The first assertion therefore applies.
\end{proof}

The same replacement theorem also gives, in every dimension,
\begin{equation}\label{eq:square-replacement}
 u(\mathcal A)\ge t(\mathcal A)^2.
\end{equation}
For its first branch this follows from $\ell_{1/d}(\varepsilon)\le\varepsilon$
and the lower bound $(1-\ell_{1/d}(\varepsilon))^2$ in its proof,
Eq.~(8); its second branch is Eq.~(11). Consequently, if
$t(\mathcal A)\ge(d+2)/(d+3)$, condition~\eqref{eq:threshold} holds:
putting $r=(d+2)/(d+3)$ gives
\[
 (d+1)r^2-(d-1)r-1=\frac{d+1}{(d+3)^2}>0,
\]
and the quadratic $(d+1)s^2-(d-1)s-1$ increases for $s\ge r$.

\begin{theorem}\label{thm:joint}
For every $d\ge2$ and $0<p<1$, $F_1^{\rm joint}(p,d)=f_d(p)$.
An optimal rank-one pair consists of the complex-space fixed-decoder
optimizer of Theorem~\ref{thm:fixed} and $\tr_S$.
\end{theorem}
\begin{proof}
The case $d=2$ is Theorem~\ref{thm:qubit}. First take $d=3$.
The product isometry $V_0\ket j=\ket{j0}$, followed by $\tr_S$, has fidelity
\begin{equation}\label{eq:product}
 F_0(q)=\frac{-128q^4+256q^3+711q^2+752q+3512}{567(q^2+8)},
 \qquad q=\sqrt{\frac{8(1-p)}{8+p}}\in(0,1).
\end{equation}
This follows by substitution in~\eqref{eq:noise}, or in the full encoder
quadratic form~\eqref{eq:quadratic}. Its exact positive gap is
\begin{equation}\label{eq:product-gap}
 F_0(q)-\frac34=
 \frac{440+3008q+1143q^2+512q^3+512q^3(1-q)}{2268(q^2+8)}>0.
\end{equation}
Take a jointly maximizing isometry $V$ and CPTP reduced recovery, whose
existence follows from compactness and the reductions above. The square
qutrit channel $\mathcal A=\T_p\C_V$ satisfies
$t(\mathcal A)\ge F_0(q)>3/4$. Lemma~\ref{lem:unitary} therefore supplies
a unitary optimal recovery $U$. Cyclicity gives
\[
 \frac19\sum_k|\tr(UM_kV)|^2
 =\frac19\sum_k|\tr(M_kVU)|^2\le f_3(p),
\]
where $M_k$ are Kraus operators of $\T_p$, and $VU$ is still an isometry.
The reverse inequality follows from feasibility of $\tr_S$.

For $d\ge4$, the same product isometry has fidelity
\[
 F_0=\lambda_3+dA-\frac{d^2-1}{d^2}C_0,
\]
where $A,C_0$ are defined in~\eqref{eq:AC}. Using $0<C_0<2A/d$,
$A=1/d-d\lambda_3$, and
$\lambda_3=2p^2/[(d^2-2)^2(d^2-1+p)]\le
2/[d^2(d^2-2)^2]$, we obtain
\begin{equation}\label{eq:product-general}
 F_0\ge m_d:=\frac{d^6-4d^4+4d^2-2}{d^4(d^2-2)}
 >\frac{d+2}{d+3}.
\end{equation}
For completeness, the numerator of the last difference, over the positive
denominator $d^4(d+3)(d^2-2)$, is
\[
 d^6-2d^5-8d^4+4d^3+12d^2-2d-6.
\]
Putting $k=d-4\ge0$ rewrites this numerator as
\[
 k^6+22k^5+192k^4+836k^3+1852k^2+1822k+434>0.
\]
The bound on $\lambda_3$ follows since $p^2/(d^2-1+p)$ increases on
$[0,1]$. Every jointly maximizing isometry thus produces a square channel
with $t(\mathcal A)\ge F_0>(d+2)/(d+3)$. Equations~\eqref{eq:square-replacement}
and~\eqref{eq:threshold} supply a unitary optimal recovery, and the same
trace cyclicity absorbs it into the encoder.
\end{proof}

\section{The exact fixed-decoder optimum in every dimension}
Put $n=d-1$ and define
\begin{equation}\label{eq:AC}
 A=\frac1d-d\lambda_3,\qquad
 C_0=\frac{2A}{d}-2(\lambda_1-\lambda_2-\lambda_3)
 \sqrt{\frac{\lambda_4\lambda_5}{d^2-1}}.
\end{equation}
These satisfy $A>0$ and $0<C_0<2A/d$. For example, with
\begin{equation}\label{eq:q}
 q=\frac ba=\sqrt{\frac{(d^2-1)(1-p)}{d^2-1+p}},\qquad
 p=\frac{(d^2-1)(1-q^2)}{d^2-1+q^2},
\end{equation}
direct simplification gives
\[
 A=\frac{(d^2-1-q^2)[d^4-4d^2+2+2(d^2-1)q^2]}
 {d(d^2-2)^2(d^2+q^2-1)},\qquad
 \frac{C_0}{2A/d}=\frac{(1-q)(d^2+q-1)}{d^2-1-q^2}\in(0,1).
\]

\begin{theorem}\label{thm:fixed}
For all $d\ge2$ and $0<p<1$, $f_d(p)$ is attained by
\begin{equation}\label{eq:family}
 V\ket0=x\ket{00}+z\sum_{i=1}^n\ket{ii},\qquad
 V\ket j=\ket{j0}\ (j>0),\qquad x\ge0,\ z\le0,\ x^2+nz^2=1.
\end{equation}
Set
\begin{gather}
 \delta=C_0(A/d-C_0/4),\quad R=A^2-AC_0+dC_0^2/4,\quad
 D(\mu)=\mu^2-A\mu+n\delta.\notag
\end{gather}
There is a unique root $A<\mu<dA$ of
\begin{equation}\label{eq:quartic}
 D(\mu)^2-n^2\bigl(R\mu^2-2nA\delta\mu+nd\delta^2\bigr)=0.
\end{equation}
For that root,
\begin{gather}
 f_d(p)=\lambda_3+\frac1d\left(An^2+\mu+
 \frac{n^2(R\mu-nA\delta)}{D(\mu)}\right),\label{eq:answer}\\
 x_* =\frac{n[(A-C_0/2)\mu-n\delta]}{D(\mu)},\qquad
 z_*=-\frac{nC_0[\mu-(2A/d-C_0/2)]}{2D(\mu)}.\label{eq:optimizer}
\end{gather}
This is also the joint optimum of Theorem~\ref{thm:joint}.
\end{theorem}

\subsection{Reduction over the full complex encoder space}
For arbitrary complex $V$, write
$r_i=\tr(D_iV)$ and $z_i=\bra\Omega V\ket i$.
Expanding the three Kraus groups of \eqref{eq:noise} yields
\begin{equation}\label{eq:quadratic}
 F(V,\tr_S)=\frac{\lambda_3}{d}\norm V_F^2+
 \frac1d\left[A\norm r^2-\sqrt d\,C_0\operatorname{Re}(r^\dagger z)
 +C_0\norm z^2\right].
\end{equation}
One way to check the expansion is to use
$C_i=QD_i^\dagger$ and
$H_i=(\ket\Omega\bra i-\sqrt{\lambda_5}C_i)/\sqrt{\lambda_4}$.
The two nonconstant Kraus contributions are
$(\lambda_1-\lambda_3)d^{-1}\sum_i|\langle C_i,V\rangle|^2$
and $\lambda_2d^{-1}\sum_i|\langle H_i,V\rangle|^2$.
Inserting $Q=aI+(b-a)P$ gives \eqref{eq:quadratic}.

The first term is $\lambda_3$ on isometries. Write the bracket as $q(V)$.
Its $2\times2$ coefficient matrix is
\[
 G=\begin{pmatrix}A&-\sqrt d C_0/2\\-\sqrt d C_0/2&C_0\end{pmatrix}\succ0.
\]
Duality for this positive quadratic form and the operator/nuclear norm
duality give
\begin{equation}\label{eq:dualnorm}
 \max_{V^\dagger V=I}q(V)=
 \max_{(\alpha,\beta)\ne0}
 \frac{\norm{I\otimes\ket\alpha+\ket\Omega\beta^T}_*^2}
 {(\alpha,\beta)^\dagger(G^{-1}\otimes I_d)(\alpha,\beta)}.
\end{equation}
Indeed the inner linear functional is
$\alpha^\dagger r+\beta^\dagger z$; its maximum absolute value over isometries is the
nuclear norm of the displayed rectangular matrix, attained by its polar
isometry. Maxima interchange after normalizing the positive denominator.

Let $a_0=\norm\alpha$, $b_0=\norm\beta$,
$c_0=\operatorname{Re}(\alpha^\dagger\beta)$, and
$k=|\alpha^\dagger\beta|^2$. The denominator in \eqref{eq:dualnorm}
depends only on $a_0,b_0,c_0$. The Gram matrix of its numerator is
\[
 a_0^2 I+\overline\beta\beta^T+
 \frac{\overline\alpha\beta^T+\overline\beta\alpha^T}{\sqrt d}.
\]
At most two eigenvalues differ from $a_0^2$. Their sum and product are
\[
 T=2a_0^2+b_0^2+2c_0/\sqrt d,\qquad
 P_2=a_0^4+a_0^2b_0^2(1-1/d)+2a_0^2c_0/\sqrt d+k/d.
\]
Consequently the nuclear norm is
$(d-2)a_0+\sqrt{T+2\sqrt{P_2}}$, which increases with $k$ at fixed
$a_0,b_0,c_0$. Collinear complex vectors with a suitable relative phase
realize $k=a_0^2b_0^2$ and preserve these other three invariants.
Zero vectors are included by continuity. Under
$M\mapsto(u\otimes\overline u)Mu^\dagger$, both vectors transform by
$\overline u$, so the collinear vectors can be placed along $\ket0$.
The polar isometry then has exactly the supports in \eqref{eq:family},
with complex amplitudes $x,z$ and a common amplitude $z$ on $\ket{ii}$,
$i>0$. Vanishing singular values admit extensions on these same supports.
Thus the complex family attains the full maximum without averaging encoders.

Substitution in \eqref{eq:quadratic} gives
\begin{align}
 q(V)={}&An^2+(A-nC_0/d)|x|^2+(n^2C_0/d)|z|^2\notag\\
 &-\frac{n(d-2)C_0}{d}\operatorname{Re}(\overline xz)
 +n(2A-C_0)\operatorname{Re}x-n^2C_0\operatorname{Re}z.\label{eq:phases}
\end{align}
Replacing $x,z$ by $|x|,-|z|$ increases or preserves every phase-dependent
term. This proves the real reduction in \eqref{eq:family}.

\subsection{A two-variable maximization and its exact dual}
Let $w=(x,z)^T$, $W=\operatorname{diag}(1,n)$ and
\[
 H=\begin{pmatrix}
 A-nC_0/d&-n(d-2)C_0/(2d)\\
 -n(d-2)C_0/(2d)&n^2C_0/d
 \end{pmatrix},\qquad
 h=\begin{pmatrix}n(A-C_0/2)\\-n^2C_0/2\end{pmatrix}.
\]
The objective is $\lambda_3+d^{-1}(An^2+w^THw+2h^Tw)$ on $w^TWw=1$.
The positive matrix $B=W^{-1/2}HW^{-1/2}$ has trace $A$ and determinant
$n\delta>0$, hence $\lambda_{\max}(B)<A$. The secular function
\[
 \norm{(\mu I-B)^{-1}W^{-1/2}h}^2
\]
strictly decreases for $\mu>A$. At $\mu=A$, its first coordinate equals
$x(A)=A(A-C_0/2)/\delta-n>1$. The last inequality follows from
\[
 A(A-C_0/2)-d\delta=A^2-3AC_0/2+dC_0^2/4>0
 \quad(0<C_0<2A/d,\ d\ge2).
\]
At $\mu=dA$, its norm is less than one because
$\norm{W^{-1/2}h}=n\sqrt R<nA$ and $\lambda_{\max}(B)<A$.
There is therefore a unique secular root in $(A,dA)$. Solving
$(\mu W-H)w=h$ gives \eqref{eq:optimizer}; imposing $w^TWw=1$ gives
\eqref{eq:quartic}, with no extraneous roots in that interval.
The formulas give $x_*>0$, $z_*<0$.

For a directly checkable certificate, set
\begin{equation}\label{eq:certificate}
 S=\mu W-H\succ0,\quad w_*=S^{-1}h,\quad
 \eta=h^TS^{-1}h=\frac{n^2(R\mu-nA\delta)}{D(\mu)},\quad
 Z=\begin{pmatrix}\eta&-h^T\\-h&S\end{pmatrix}\succeq0.
\end{equation}
The Schur complement is zero, and
$Z(1,w_*^T)^T=0$. On $w^TWw=1$ the exact identity
\begin{equation}\label{eq:gap}
 f_d(p)-F(V(w),\tr_S)=d^{-1}(w-w_*)^TS(w-w_*)\ge0
\end{equation}
proves the bound with equality at $w_*$. Equivalently $Z$ is dual to
the lifted quadratic program and $(1,w_*^T)^T(1,w_*^T)$ is a matching
rank-one primal matrix. This completes Theorem~\ref{thm:fixed}.

\section{A decoder certificate for every encoder in the family}
For a fixed isometry $V$, put $\mathcal A=\T_p\C_V$ and
$K=J(\mathcal A^\dagger)/d^2$, where the input-first Choi convention is
$J(\R)=\sum_{ij}\ket i\bra j\otimes\R(\ket i\bra j)$.
The decoder SDP and its dual are
\begin{equation}\label{eq:decoder-sdp}
 \max_{X\succeq0,\ \tr_{\rm out}X=I_d}\tr(KX),\qquad
 \min_{Y=Y^\dagger,\ Y\otimes I_d\succeq K}\tr Y.
\end{equation}
These normalizations follow immediately from \eqref{eq:fidelity}.

\begin{proposition}\label{prop:decoder}
For every real encoder \eqref{eq:family}, the identity reduced recovery,
and therefore the physical decoder $\tr_S$, is optimal among all decoders.
\end{proposition}
\begin{proof}
Write $z=-t$, $t\ge0$, $s=x-nt$, and put
\[
 \alpha=d\lambda_3,\quad\beta=dA,\quad c=d^2C_0/2\in(0,\beta),
\]
\[
 a_0=\alpha+\beta x^2-c(2sx/d-2s^2/d^2),\quad
 b_0=\alpha+\beta t^2+c(2st/d+2s^2/d^2),\quad
 g=\beta x-cs/d.
\]
With $P_0=\ket0\bra0$, $\Pi=I-P_0$, direct evaluation of
\eqref{eq:factorization} gives
\[
 \mathcal A(\ket0\bra0)=a_0P_0+b_0\Pi,\quad
 \mathcal A(\ket j\bra j)=\alpha I+\beta\ket j\bra j\ (j>0),
\]
\[
 \mathcal A(\ket0\bra j)=g\ket0\bra j\ (j>0),\qquad
 \mathcal A(\ket i\bra j)=\beta\ket i\bra j\ (i,j>0,\ i\ne j).
\]
The adjoint relations determine the remaining matrix units. Take
\begin{equation}\label{eq:Y}
 Y=d^{-2}\big[(a_0+ng)P_0+(g+n\beta+\alpha)\Pi\big].
\end{equation}
On the span of $\ket{00}$ and $n^{-1/2}\sum_{j>0}\ket{jj}$,
and on its orthogonal invariant subspaces, the slack is
\begin{equation}\label{eq:decoderblocks}
 \begin{aligned}
 d^2(Y\otimes I-K)\ &\cong
 g\begin{pmatrix}n&-\sqrt n\\-\sqrt n&1\end{pmatrix}
 \oplus(a_0+ng-\alpha)I_n\\
 &\quad\oplus(g+n\beta+\alpha-b_0)I_n
 \oplus(g+n\beta)I_{n^2-1}.
 \end{aligned}
\end{equation}
The last block is absent when $d=2$. To see positivity, first
$g=(\beta-c/d)x+cnt/d\ge0$. Also
$s(x+t)\le x^2$, since $x,t\ge0$. Therefore
\[
 a_0-\alpha=\beta x^2-\frac{2cn}{d^2}s(x+t)
 \ge(\beta-2cn/d^2)x^2\ge0,
\]
\[
 b_0-\alpha=\beta t^2+\frac{2c}{d^2}s(x+t)
 \le\beta t^2+\frac{\beta}{n}x^2=\frac\beta n.
\]
Here $2cn/d^2\le\beta$ follows from $c<\beta$ and $2n\le d^2$.
Every block in \eqref{eq:decoderblocks} is positive semidefinite.
The slack annihilates $\ket I\!\rangle=\sum_i\ket{ii}$,
the Choi vector of the identity recovery. Its dual value thus equals
the identity recovery's primal value, proving optimality.
\end{proof}

\section{Exact instances and a relaxation obstruction}
For $d=2$, the formula is particularly simple. Define
\[
 k=\frac{3-2p}{3+p},\quad h_0=\frac{1+p}{2},\quad
 w_0=\sqrt{\frac{(3+p)(1-p)}{12}},\quad
 A_2=k(h_0+w_0),\quad B_2=k(h_0-w_0).
\]
If $\xi>A_2$ is the unique root there of
\[
 \xi^4-2(A_2+B_2)\xi^3+4A_2B_2\xi^2-(A_2B_2)^2=0,
\]
then the joint fidelity is
\[
 F_1^{\rm joint}(p,2)=\frac14\left(1+\xi+
 \frac{A_2^2}{\xi-A_2}+\frac{B_2^2}{\xi-B_2}\right),\quad
 x_* =\frac{A_2}{\xi-A_2},\quad z_*=-\frac{B_2}{\xi-B_2}.
\]
This is also obtained directly from the circle objective
$[1+A_2(x^2+2x)+B_2(z^2-2z)]/4$.
At $p=9/13$, $A_2=105/208$, $B_2=49/208$, and
$F_1^{\rm joint}\simeq0.6462632754$.

For a qutrit instance choose $q=1/2$, hence $p=8/11$. Then
\[
 \lambda_3=\frac4{1617},\quad A=\frac{527}{1617},\quad
 C_0=\frac{578}{4851}.
\]
Writing $T=1617\mu$, the specified quartic becomes
\begin{equation}\label{eq:qutritquartic}
 T^4-1054T^3-440147T^2+155349060T-12068784500=0,
 \qquad 1305<T<1306.
\end{equation}
Equations~\eqref{eq:optimizer}, \eqref{eq:certificate}, and \eqref{eq:Y}
give exact encoder and decoder certificates over this real root field.
Numerically $x_*\simeq0.9632339434$, $z_*\simeq-0.1899741696$, and
\[
 F_1^{\rm joint}(8/11,3)=f_3(8/11)\simeq0.8922491436
 <F^{\rm unrestricted}(8/11,3)=\frac{69}{77}.
\]
In particular, Theorem~\ref{thm:joint} proves the global certificate
\begin{equation}\label{eq:open}
 \tr(K(V)X)\le f_3(8/11)\quad
 \text{for all }V^\dagger V=I_3,\ X\succeq0,\ \tr_{\rm out}X=I_3.
\end{equation}
The root isolation, fixed-decoder gap matrix and recovery criterion together
certify this value over the entire complex joint feasible set.

There is an exact obstruction to the simplest polynomial certificate even
with the decoder fixed. No identity expressing $f_d(p)-F(V,\tr_S)$ as a sum
of squares of affine real polynomials in $V$, plus constant real multipliers
of the equations $V^\dagger V-I=0$, can hold. Indeed let
$C_i=QD_i^\dagger$, $i=0,\ldots,d-1$, be the unrestricted optimal encoder's
Kraus operators~\cite[Eq.~(19)]{LiJiang}. Assign probability $1/(2d)$ to
each $V=\pm\sqrt d\,C_i$. This distribution is not supported on isometries,
but it satisfies
\[
 \mathbb E[V^\dagger V]=I,\qquad
 \mathbb E[F(V,\tr_S)]=\sum_i F(C_i,\tr_S)=\lambda_1>f_d(p).
\]
The last strict inequality follows from the rank requirement for unrestricted
optimality in Ref.~\cite{LiJiang}, Sec.~III. Taking expectations in the
hypothetical identity is a contradiction. Thus a quadratic moment relaxation
loses precisely the encoder constraint at issue. This obstruction does not
rule out degree-four certificates. The qutrit proof above bypasses this
relaxation by applying the recovery criterion before optimizing the encoder.

\section{A uniform question beyond this family}
Equation~(16) of Ref.~\cite{LiJiang} already bounds the gap with both encoder
and decoder optimized. This family approaches its small-error ceiling to
second order as $d$ grows; it does not attain that bound at finite $d$.
A precise broader question is the sharp coefficient at fixed logical and
physical dimensions. For qubits encoded into dimension four, let
\[
 c_{2,4}=\limsup_{\eta\downarrow0}\ \sup_{\substack{
 \mathcal N:\mathsf L(\mathbb C^4)\to\mathsf L(\mathbb C^4)\ {\rm CPTP}\\
 0<\varepsilon^{\rm opt}(\mathcal N)\le\eta}}
 \frac{\Delta^{\rm opt}(\mathcal N)}{[\varepsilon^{\rm opt}(\mathcal N)]^2},
\]
where $\varepsilon^{\rm opt}=1-F^{\rm unrestricted}$ and
$\Delta^{\rm opt}=F^{\rm unrestricted}-F_1^{\rm joint}$ at $\tau=I_2/2$.
Their Eqs.~(16) and (25) give
\[
 \frac19\le c_{2,4}\le\frac14.
\]
Resolving this interval requires either a new noise family approaching
$1/4$ or a quantitative incompatibility among the near-equality conditions
in the rank-one replacement bound and joint optimality. This is distinct
from finding an exact solution for the present model.

\paragraph{Verification and computer assistance.}
The conclusions above are mathematical proofs, not Lean formalizations.
The accompanying source checks symbolic identities and finite exact matrix
certificates with rational arithmetic; numerical approximations serve to
identify values. Haobo Ma and Wenlin Zhang used AI tools for search,
double-checking and proofreading.

\begin{thebibliography}{9}
\bibitem{LiJiang} B. Li and L. Jiang,
\emph{High-Rank Encoding Can Improve Approximate Quantum Error Correction},
arXiv:2609.00778v1 (2026),
\url{https://arxiv.org/abs/2609.00778v1}.
\bibitem{KBF} V. Karimipour, F. Benatti, and R. Floreanini,
\emph{Quasi inversion of qubit channels}, Phys. Rev. A \textbf{101},
032109 (2020), doi:10.1103/PhysRevA.101.032109,
\url{https://arxiv.org/abs/1909.06118v5}.
\bibitem{SSMZK} F. Shahbeigi, K. Sadri, M. Moradi, K. \.Zyczkowski,
and V. Karimipour,
\emph{Quasi-inversion of quantum and classical channels in finite dimensions},
J. Phys. A: Math. Theor. \textbf{54}, 345301 (2021),
doi:10.1088/1751-8121/ac13db,
\url{https://arxiv.org/abs/2104.06062v3}.
\end{thebibliography}
\end{document}
